\documentclass[a4paper,11pt]{article}
\usepackage[T1]{fontenc}
\usepackage[polish]{babel}
\usepackage{lmodern}
\usepackage[margin=2.35cm]{geometry}
\usepackage{microtype}
\setlength{\parindent}{1.2em}
\setlength{\parskip}{0.25em}

\begin{document}
\begin{center}
  {\Large\bfseries Paczka o siódmej}\\[0.3em]
  {\small Próba pastiszu dziewiętnastowiecznej prozy realistycznej}
\end{center}

O siódmej rano, o której ludzie sumienni myślą już o obowiązkach,
a ludzie mniej sumienni o tym, jak je odłożyć, panna Lena ujrzała
pod drzwiami paczkę wielkości porządnego słownika. Nie było na niej
nazwiska nadawcy. Ktoś wypisał tylko: „Otworzyć przed ósmą”.
Na klatce schodowej nie znalazła żywej duszy; przez okno deszcz
ukazywał podwórze, które nawet w pogodę nie uchodziło za piękne.
Lena położyła karton na stole i spojrzała na zegar. Zegar, jak
wszystkie zegary w chwilach niepewności, pracował bez współczucia.

W środku leżał klucz, kartka z numerem 12 i mapa osiedla. Czerwony
ołówek wskazywał ławkę przy nieczynnej fontannie. Ileż razy Lena
mijała to miejsce, nie przypuszczając, że może służyć czemuś więcej
niż odpoczynkowi gołębi! Zadzwoniła do sąsiada naprzeciwko, lecz
telefon milczał. Zrobiła więc zdjęcie paczki, schowała kartkę i
wyszła, bo w mieście człowiek rozsądny nie czeka bez końca na
wyjaśnienie zagadki, jeżeli zagadka wyznacza godzinę.

Przy fontannie siedział pan w szarym płaszczu, pochłonięty gazetą
z miną właściwą ludziom, którzy wiedzą, co zaszło w całym świecie,
a nie widzą niczego pod własnymi nogami. Lena przeszła obok dwa razy.
Dopiero za trzecim spojrzeniem dostrzegła skrzynkę pod ławką.
Zamek opierał się chwilę, po czym ustąpił kluczowi z paczki.
W skrzynce znajdowały się notes, ołówek i wiadomość: „Zanieś notes
do biblioteki przed ósmą. Archiwum zamykają dziś na remont”.

Była siódma trzydzieści osiem. Lena ruszyła ku bibliotece, ale pod
daszkiem sklepu otworzyła notes. Miękka okładka i zagięte rogi nie
zapowiadały skarbu. W środku znalazła jednak szkice starych drzew,
nazwy ulic i daty, z których układała się pamięć o osiedlu;
taką pamięć najłatwiej stracić wtedy, gdy wszyscy są przekonani,
że zachowa ją ktoś inny.

Do biblioteki dotarła przed ósmą. Bibliotekarka rozpoznała pismo
dawnego mieszkańca: właśnie tych zapisków brakowało do wystawy
o dzielnicy. Lena oddała notes, poprosiła o potwierdzenie i wróciła
do domu. Pusty karton czekał na stole. Na odwrocie klapy, dotąd
niewidocznym, widniał podpis: „Dziękuję za pomoc -- sąsiad spod
dwunastki”. Dziwne było tylko to, że najbliższa droga do rozmowy
z sąsiadem prowadziła przez całe osiedle.

\end{document}
